\subsection{AI, ML, and Deep Learning}
\begin{frame}{\insertsubsection}
	\centering\pic[width=.5\linewidth]{ml/AI_relations} %TODO build with tikz
\end{frame}


\subsection{Artificial Intelligence}
\begin{frame}{\insertsubsection{}  \mytitlesource{\href{https://www.pearson.de/artificial-intelligence-a-modern-approach-global-edition-9781292401133}{Russel-Norvig}}}
	\begin{fancycolumns}[animation=none]
		\begin{definition}{Artificial Intelligence}
			Area of study that develops methods to allow machines to perceive their environment and take actions that maximize the chance of achieving defined goals
		\end{definition} \pause
		\begin{note}{Common Tasks}
			\begin{itemize}
				\item Reasoning
				\item Knowledge Representation
				\item Planning
				\item Natural Language Processing
				\item Perception
				\item Support for robotics
				\item \dots
			\end{itemize}
		\end{note} \pause
		\nextcolumn
		\begin{definition}{Automated Planning}
			Automated creation of strategies or action sequences to achieve a predefined goal
		\end{definition} \pause
		\begin{example}{Keywords}
			\begin{itemize}
				\item State: collection of currently active fluents
				\item Action: changes state
				\item Precondition: fluents necessary to apply action
				\item Effect: changes to state after action is applied
			\end{itemize}
		\end{example} \pause
		\begin{note}{Solution Strategies}
			\begin{itemize}
				\item Forward state-space search
				\item Backward relevant-states search
				\item Heuristic search
			\end{itemize}
		\end{note} 
	\end{fancycolumns}
\end{frame}

\subsection{Machine Learning}

\xkcdframe{1838}

\begin{frame}{\insertsubsection{} \mytitlesource{\href{https://www.statlearning.com/}{Statistical Learning}}}
	\begin{fancycolumns}[animation=none]
		\onslide<1->{
		\begin{definition}{Machine Learning}
			Area of research of estimating an unknown target function from multiple predictors based on example data
		\end{definition}}
		\onslide<3->{
		\begin{example}{Common Models}
			\small
			\begin{itemize}
				\item Decision trees (often with gradient-boosting or bootstrap aggregating)
				\item Support vector machines
				\item Regression
				\item Bayesian networks
				\item Gaussian processes
				\item Rule-based models
				\item \dots
			\end{itemize}
		\end{example}}
		\nextcolumn
		\onslide<2->{\begin{note}{Common Steps}
				\small
			\begin{enumerate}
				\item Feature extraction
				\item Feature transformation
				\item Feature selection
				\item ML model training
			\end{enumerate}
		\end{note}} 
		\onslide<4->{\begin{example}{Support Vector Machines}
			\centering\pic[width=.6\columnwidth]{ml/svm}
		\end{example}}
	\end{fancycolumns}
\end{frame}

\subsection{Neural Networks}
\begin{frame}{\insertsubsection{} \mytitlesource{\href{https://www.statlearning.com/}{Statistical Learning}}}
	\begin{fancycolumns}[animation=none]
		\begin{definition}{Perceptron}
			$$f(\mathbf{x}) = h(\mathbf{w} \cdot \mathbf{x} + b)$$
			Weights $\mathbf{w}$, input $\mathbf{x}$, bias $b$
		\end{definition} \pause
		\begin{definition}{Neuron}
			$$y_k = \varphi \left(\sum_{j=0}^m w_{kj} x_j \right)$$
			Activation function $\varphi$, bias $w_{k0}$
		\end{definition} \pause
		\begin{definition}{Neural Network}
			System of layered neurons with varying interconnections, activation functions, and weights
		\end{definition} \pause
		\nextcolumn
		\begin{example}{Exemplary Neural Network}
			\centering\pic[width=.75\columnwidth]{ml/neural_network}
		\end{example}
	\end{fancycolumns}
\end{frame}

\subsection{Large Language Models}
\begin{frame}{\insertsubsection{} \mytitlesource{\href{https://web.stanford.edu/~jurafsky/slp3/ed3book_jan72023.pdf}{Introduction to NLP}}}
	\begin{fancycolumns}[animation=none]
		\begin{definition}{Large Language Models}
			Neural network specialized on language processing
			\vspace{-1em}
			\begin{itemize}
				\item vast textual training data set
				\item word-wise output prediction
				\item uses transformer architecture
			\end{itemize} 
		\end{definition} \pause
		\begin{note}{Transformer Architecture}
			\begin{enumerate}
				\item Tokenization
				\item Embedding
				\item Positional Encoding
				\item Transformer Blocks (Attention)
				\item Softmax-Layer
			\end{enumerate}
		\end{note} \pause
		\nextcolumn
		\begin{example}{Transformer Architecture}
			\centering\pic[width=.8\columnwidth]{ml/full_transformer}
		\end{example}
	\end{fancycolumns}
\end{frame}

\subsection{Agents}
\begin{frame}{\insertsubsection{} \mytitlesource{\href{https://www.statlearning.com/}{Statistical Learning}}}
	\begin{fancycolumns}[animation=none]
		\begin{definition}{Traditional AI Agent}
			\begin{itemize}
				\item Perceives environment through sensors
				\item Acts upon environment through actors
				\item Tries to achieve a given goal
			\end{itemize}
		\end{definition} \pause
		\begin{definition}{LLM-based Agents}
			\begin{itemize}
				\item Uses LLM as core reasoning engine
				\item Additional perception, memory, decision-making, and action modules
				\item Perception - Decision - Execution
				\item Multi-agent systems
			\end{itemize}
		\end{definition} \pause
		\nextcolumn
		\begin{example}{Traditional Reflex Agent}
			\centering\pic[width=\columnwidth]{ml/trad_agent}
		\end{example}
	\end{fancycolumns}
\end{frame}
